\section{Arquitetura}

\subsection{Mecanismo anti-spoiler}

\subsection{Validação}

Os modelos geradores decidem \textit{o que} o assistente responde a cada solicitação do usuário e o mecanismo anti-spoiler descrito na Seção~\ref{subsec:antispoiler} garante que o assistente nunca exponha o que o leitor ainda não alcançou; nenhum deles, porém, estabelece se uma resposta gerada está de fato \textit{correta}, nem comunica ao leitor o quanto o sistema confia nela. É essa lacuna que a camada de validação preenche. Atribuí-la ao sistema não é um detalhe de implementação: a ``comunicação cuidadosa das respostas, exibindo limitações e incerteza'' figura entre os requisitos \textit{Must} do projeto (Tabela~1), e esta seção descreve o mecanismo --- o \textit{validador} --- com que o realizamos.

O ponto de partida da camada de validação é o reconhecimento de que ``correto'' não é uma propriedade única. Uma resposta gerada pelo assistente costuma reunir afirmações de naturezas distintas: um fato de enredo ancorado no texto já lido, uma paráfrase que precisa preservar o sentido do trecho original, um sentido de palavra que deve corresponder ao registrado em dicionário, ou ainda um fato sobre o mundo real que vive fora da obra. Cada uma dessas naturezas se verifica de forma diferente e, sobretudo, contra uma fonte de evidência diferente. Submeter todas a um único \textit{prompt} do tipo ``isto está correto?'' herdaria justamente o pior comportamento desse tipo de instrução: o erro confiante, em que o modelo afirma com a mesma segurança um fato que recuperou do texto e um que apenas recordou de memória --- uma alucinação. É essa indistinção que a validação precisa desfazer, tratando cada tipo de afirmação com a metodologia de checagem mais apropriada a ele.

Essa correspondência entre tipo de afirmação e checagem se concretiza de modo diferente conforme a intenção que o leitor aciona (ver Seção 7.1). Para contextualizar e relembrar, o validador decompõe a resposta gerada em afirmações menores e auto-contidas --- resolvendo pronomes e separando fatos de atribuições causais --- e verifica cada uma isoladamente contra o contexto recuperado, de modo que cada afirmação receba seu próprio veredito. A definição segue uma estrutura própria: em vez de decompor a resposta, o validador checa a definição geral do trecho como uma unidade e, em \textit{prompts} separados, a definição de cada palavra difícil sinalizada (Seção~\ref{subsec:dificuldade-palavras}), apoiando esta última etapa em um dicionário offline (WordNet, via \textit{nltk}) que serve de fonte de evidência lexical. A paráfrase, por fim, é entregue ao validador de forma integral, sem decomposição: o que se verifica é se a frase reescrita preserva o sentido do trecho original, uma relação de equivalência que se perderia caso a frase fosse fragmentada.

Um segundo princípio orienta o fluxo de validação: um simples \textit{correto/incorreto} nos dá pouca informação sobre a resposta gerada. Uma afirmação pode estar parcialmente sustentada pela evidência, ou ser plausível sem que se tenha como confirmá-la apenas com o contexto fornecido, e tratar todos esses casos como um único ``incorreto'' apagaria distinções que importam ao leitor. Por isso o validador produz duas saídas distintas: um \textit{veredito} graduado, que estima o quanto cada afirmação se sustenta na evidência consultada, e uma estimativa de incerteza sobre esse veredito, que captura o quanto o validador "confia" no julgamento que acabou de emitir \cite{liu2025uncertainty,kadavath2022language}. Essa estrutura permite à interface fazer ressalvas transparentes -- sinalizar que uma resposta é plausível mas não pôde ser confirmada, em vez de forçá-la num binário que ocultaria essa diferença. 

Por fim, a validação não é um \textit{prompt} genérico único, mas um conjunto de módulos especializados por tipo de afirmação: cada checagem torna-se mais simples e mais alinhada àquilo que de fato precisa verificar, possuindo um \textit{prompt} próprio elaborado em torno de objetivos específicos. As subseções a seguir detalham, nessa ordem, como o validador está integrado na arquitetura e suas fontes de verificação, o \textit{pipeline} completo de validação, a comunicação de incerteza e os estados da interface e, por fim, a avaliação \textit{offline} do mecanismo como método de calibração e seus resultados.

\subsubsection*{Integração e fontes de ancoragem}
\label{subsubsec:arquitetura}

O validador é a última camada de todos os caminhos de requisição: qualquer que seja a intenção acionada pelo leitor --- definir, parafrasear, contextualizar ou relembrar (ver Seção~\ref{visao-geral-sistema}) --, a resposta gerada só chega à interface depois de passar por ele, e é implementado como um modelo distinto daquele que gera as respostas. O ponto central de seu encaixe na arquitetura é a fonte sobre a qual ele se apoia para julgar: o validador ancora cada checagem no mesmo material que o gerador correspondente viu -- os trechos recuperados dentro do limite de posição do leitor, no caso da contextualização e do relembrar; o contexto reduzido em torno da passagem, no caso da definição; o próprio trecho selecionado, no caso da paráfrase. Como esse material já passou pelo filtro determinístico de posição da recuperação (Seção~\ref{subsec:antispoiler}), o julgamento do validador herda a mesma abordagem anti-spoiler da geração: ele não tem acesso a nada além do que o leitor já leu, e portanto não deve validar uma resposta recorrendo a texto que o próprio leitor ainda não alcançou. Vale notar que ancorar a checagem nos trechos recuperados, e não em toda a obra, é também uma decisão de qualidade: inundar o validador com dezenas de capítulos por afirmação seria custoso e degradaria a precisão da verificação de implicação textual, ao oferecer muito mais texto irrelevante capaz de induzir um veredito equivocado.
 
Chamamos aqui de \textit{afirmação} cada unidade verificável que compõe uma resposta gerada. Em alguns casos, a resposta é decomposta em várias afirmações menores e auto-contidas, cada uma julgada isoladamente: é o que ocorre na contextualização, no relembrar e também na definição de um trecho, em que retornamos a explicação geral do trecho selecionado. Em outros casos, a unidade a verificar é a própria resposta, tomada de forma integral --- a paráfrase inteira e a definição de cada palavra difícil encontrada no trecho em que o usuário requisitou definição. A natureza da afirmação determina não só como ela é checada, mas contra que fonte. Além dos trechos recuperados, o validador recorre a um dicionário \textit{offline} (WordNet) para a validação léxica das definições, confrontando o sentido proposto com os sentidos ali registrados. Além disso, uma afirmação pode ser classificada como conhecimento de mundo -- fatos históricos, lugares reais, pano de fundo cultural que vivem fora da obra. Nesse caso, a verificação exigiria uma fonte externa, consultada através de uma busca na web. Trata-se do caso mais difícil, e optamos por deixá-lo fora desta prova de conceito: uma checagem factual robusta contra a web traz desafios próprios (formulação da consulta, qualidade e confiabilidade das fontes, distinção entre ficção e mundo real) que extrapolam o nível de maturidade TRL 3 a que nos propusemos. Na prática, afirmações de conhecimento de mundo são tratadas conservadoramente como não verificáveis, conforme detalhado em seções posteriores. A Tabela~\ref{tab:fontes-verificacao} sintetiza, para cada tipo de afirmação, a fonte de verificação correspondente e critério de validação.


\begin{table}[ht]
  \centering
  \caption{Tipos de afirmação, suas fontes de verificação e o critério aplicado a cada uma.}
  \label{tab:fontes-verificacao}
  \begin{tabularx}{\textwidth}{>{\raggedright\arraybackslash}p{0.22\textwidth} >{\raggedright\arraybackslash}X >{\raggedright\arraybackslash}X}
  \toprule
  \textbf{Tipo de afirmação} & \textbf{Fonte de verificação} & \textbf{Critério de validação} \\
  \midrule
  Ancorada no contexto (identidade de personagens, fatos do enredo, sentido de um trecho) & O texto lido até a posição atual do leitor, recuperado dentro do limite anti-spoiler & A afirmação precisa ser sustentada pelo trecho recuperado, por implicação textual, sem recorrer a informações que o leitor ainda não alcançou. \\
  \addlinespace
  Equivalência de sentido (paráfrase) & O trecho originalmente selecionado pelo leitor & A paráfrase precisa preservar o sentido do trecho de forma bidirecional, sem acrescentar, omitir nem distorcer o conteúdo. \\
  \addlinespace
  Léxico-semântica (definição) & Um dicionário \textit{offline} (WordNet, via \textit{nltk}) e o pequeno contexto recuperado ao redor da passagem, também sujeito ao filtro anti-spoiler & Dupla exigência: a definição precisa ser válida no dicionário e, ao mesmo tempo, corresponder ao sentido adequado àquela passagem. \\
  \addlinespace
  Conhecimento de mundo (história, lugares reais, pano de fundo) & Fontes externas, por meio de busca na web --- recurso não implementado nesta prova de conceito & Checagem de factualidade contra a evidência externa recuperada; o caso mais difícil e, por ora, fora do escopo. \\
  \bottomrule
  \end{tabularx}
  \end{table}




\subsubsection*{O fluxo de validação}
\label{subsubsec:pipeline}

Toda resposta percorre um fluxo cuja espinha dorsal é compartilhada, ainda que o trecho inicial varie por funcionalidade, como sintetiza a Figura~\ref{fig:subpipeline-validacao}. Os caminhos de contextualização e de relembrar passam pelas três etapas iniciais --- decompor, rotear e emitir um veredito por afirmação; a definição reaproveita esse mesmo caminho para o sentido geral do trecho e o combina com uma checagem lexical à parte; a paráfrase salta direto para um veredito único sobre a resposta inteira. A partir daí, todos convergem: os vereditos individuais são agregados em um veredito final, que passa por um filtro de incerteza antes de chegar ao leitor.

A decomposição, presente apenas nos caminhos que a utilizam, é feita por um modelo LLM (Claude Sonnet) encarregada unicamente de quebrar a resposta em afirmações atômicas --- resolvendo pronomes e separando cada fato de eventuais atribuições causais a ele atreladas, separação que se mostrou necessária para evitar vereditos parciais equivocados --- e de classificar cada afirmação em um de dois tipos: \texttt{contexto} ou \texttt{conhecimento de mundo}. Afirmações de conhecimento de mundo sofrem um curto-circuito imediato para o rótulo \textit{Não verificável}, sem nova chamada ao modelo, já que a busca na web que as confirmaria não foi implementada (ver Seção~\ref{subsubsec:arquitetura}); as demais seguem para o veredito.

O veredito por afirmação é uma checagem de implicação textual contra o contexto recuperado, feita por um modelo LLM (Claude Sonnet) a partir de um \textit{prompt} próprio, que retorna um entre quatro rótulos --- \textit{Sustentada}, \textit{Parcialmente sustentada}, \textit{Contradita} e \textit{Não verificável} --- acompanhado de uma confiança verbalizada e de uma justificativa de uma frase. \textit{Não verificável} é o rótulo atribuído a afirmações que o validador não consegue ancorar na evidência fornecida, e não uma falha de execução; ele mantém o validador ``honesto'' diante de uma afirmação plausível mas sem suporte no que o leitor já leu. Os vereditos das afirmações de uma mesma resposta são então agregados pelo \textit{pior caso}: o veredito final é o rótulo mais severo entre eles, de modo que uma única afirmação contradita basta para impedir que a resposta seja apresentada como confiável --- um padrão deliberadamente conservador, adequado a um auxílio à leitura. A confiança verbalizada associada a esse veredito final é a da própria afirmação de pior caso, que determinou a agregação.

O veredito por definição segue a mesma lógica, aplicada a duas fontes em paralelo. Cada palavra difícil é validada em um \textit{prompt} separado, no qual o validador confronta o par palavra--definição com a entrada correspondente do dicionário \textit{offline} WordNet (quando existente) e com o contexto recuperado; o sentido geral do trecho, por sua vez, é validado pelo caminho ancorado em contexto descrito acima. Cada uma dessas checagens devolve o mesmo conjunto de saídas do veredito por afirmação (rótulo, confiança e justificativa), e o veredito final é obtido pela mesma agregação pelo pior caso.

O veredito por paráfrase dispensa decomposição e agregação: a frase parafraseada é checada em passe único e bidirecional contra o trecho selecionado, verificando que nada foi acrescentado, omitido ou distorcido. A checagem devolve as mesmas saídas dos demais vereditos, e o rótulo resultante já é, em si, o veredito final.

Qualquer que seja o caminho, o veredito final passa por uma etapa que chamamos de filtro de incerteza. Ela aplica um método de calibração, implementado \textit{offline} e detalhado na seção a seguir, que estima o quanto se pode de fato confiar na confiança verbalizada pelo modelo. Quando a incerteza estimada é baixa, o veredito é aceito e apresentado como tal; quando é alta, o veredito é rejeitado e o leitor é avisado de que a resposta não pôde ser confirmada com segurança.

Esse resultado é comunicado na interface por meio de três informações. A primeira é o estado final, em um de três rótulos: \textit{Válido} e \textit{Não confiável} aparecem quando o filtro de incerteza considera o veredito digno de confiança (um para a resposta sustentada, outro para a contradita), ao passo que \textit{Com ressalvas} aparece quando há incerteza alta demais sobre o próprio veredito, ou quando o veredito final considerou o resultado como \textit{Não verificável}. A segunda informação é o detalhamento de cada veredito intermediário --- por afirmação, por definição ou por paráfrase ---, com o rótulo atribuído e a justificativa que o modelo deu. A terceira são as fontes de verificação efetivamente consultadas ao longo do processo. Em conjunto, as três compõem uma trilha auditável: em vez do raciocínio bruto do validador, o leitor recebe a evidência usada, o que foi checado e por que se chegou àquele veredito.



\subsubsection*{Comunicação de incerteza}
\label{subsubsec:incerteza}

Cada veredito vem acompanhado de uma confiança verbalizada: um número entre 0 e 1 que o próprio modelo reporta ao lado do rótulo, dizendo o quanto se julga seguro. O problema é que esse autorrelato é notoriamente superconfiante --- modelos de linguagem tendem a declarar confiança alta mesmo quando erram --, e tomá-lo ao pé da letra esvaziaria o sentido de comunicar incerteza ao leitor. O filtro de incerteza existe para corrigir isso: em vez de confiar no número reportado pelo modelo, ele estima o quanto esse número de fato prediz acerto, e só então decide se o veredito pode ser apresentado como confiável.

Uma dificuldade prática moldou essa estimativa. A API do modelo utilizado não expõe as probabilidades dos tokens de saída (\textit{logprobs}), de modo que não temos acesso à incerteza interna da geração; restam apenas sinais de ``caixa-preta'', medidos sobre as respostas observáveis. Nossa primeira tentativa foi medir a \textit{consistência intra-execução}: rodar o mesmo \textit{prompt} um número $N$ de vezes e observar o quanto o veredito variava entre as rodadas \cite{Huang2025lookbefore}. O sinal mostrou-se inerte --- quando o modelo está confiante sobre uma entrada, ele repete o mesmo veredito em todas as rodadas, e a variação não aparece mesmo nos casos em que o veredito está errado. Reorientamos então o esforço para uma questão diferente: em vez de variar a mesma entrada, perguntar se, \textit{ao longo de muitas entradas distintas}, a confiança verbalizada de fato separa os vereditos certos dos errados. Essa é uma pergunta de calibração, e respondê-la exige uma referência externa.

A referência é um conjunto de \textit{anotações-ouro} (\textit{gold-labels}): um conjunto de respostas cujos vereditos corretos foram atribuídos manualmente, contra o qual medimos o desempenho do validador (detalhado na seção a seguir). Sobre ele adotamos a estratégia de \textit{predição seletiva} (\cite{yaniv2010selective,geifman2017selective}): mantemos a confiança verbalizada do modelo como \textit{score} e calibramos um limiar $\tau$ de ``compromisso''. O validador compromete-se com o veredito quando a confiança fica acima de $\tau$ --- exibindo-o como \textit{Válido} ou \textit{Não confiável} --- e se abstém quando fica abaixo, recolhendo a resposta ao estado intermediário \textit{Com ressalvas} (\cite{yadkori2024mitigating}). O limiar não é arbitrário: ele é lido da curva de risco--cobertura (\textit{risk--coverage}) do conjunto de anotações-ouro, no ponto que entrega a acurácia desejada entre os vereditos comprometidos, e só depois de um teste que confirme que a confiança verbalizada é, naquele conjunto, discriminativa o bastante para valer como \textit{score} de abstenção.

A Tabela~\ref{tab:estados-ui} resume como a saída do validador, depois do filtro, se traduz nos três estados apresentados ao leitor.

\begin{table}[ht]
\centering
\caption{Mapeamento da saída do validador para os três estados da interface.}
\label{tab:estados-ui}
\begin{tabularx}{\textwidth}{>{\raggedright\arraybackslash}p{0.30\textwidth} >{\raggedright\arraybackslash}p{0.20\textwidth} >{\raggedright\arraybackslash}X}
\toprule
\textbf{Saída do validador} & \textbf{Estado na interface} & \textbf{Comportamento} \\
\midrule
Sustentada, com baixa incerteza & \textit{Válido} & A resposta é exibida com um selo de validação e a trilha auditável fica acessível. \\
\addlinespace
Contradita ou sem suporte, com baixa incerteza & \textit{Não confiável} & O leitor é avisado de que não foi possível dar uma resposta confiável, com a opção de tentar novamente. \\
\addlinespace
Alta incerteza, ou veredito parcialmente sustentado ou não verificável & \textit{Com ressalvas} & A resposta é exibida com a marca de que não pôde ser inteiramente verificada, expondo a trilha e as fontes consultadas. \\
\bottomrule
\end{tabularx}
\end{table}

A existência do estado \textit{Com ressalvas} é o ponto central deste desenho. Uma resposta pode estar correta e, ainda assim, não ser verificável com a evidência disponível --- caso típico das afirmações de conhecimento de mundo, que o validador não tem como confirmar. Colapsar esses casos em \textit{Não confiável} seria impreciso: trataria como suspeita uma resposta que apenas não pôde ser checada, e, com o tempo, ensinaria o leitor a desconfiar justamente da franqueza do sistema. Manter um estado separado para a incerteza é o que permite ao assistente admitir os limites do que sabe sem se autodepreciar a cada dúvida.


\subsubsection*{Calibração \textit{offline}}
\label{subsubsec:calibracao}

O limiar $\tau$ de ``compromisso'' mencionado anteriormente não é escolhido a priori: ele é o produto de um experimento de calibração conduzido inteiramente \textit{offline}, separado do caminho de execução --- roda uma vez, sobre dados de referência, e não a cada resposta ao leitor. É esse experimento que sustenta nossa alegação de maturidade TRL 3 para o validador do sistema: ele é a evidência de que o validador se comporta como descrevemos.

Conforme mencionado anteriormente, o experimento parte de um conjunto de \textit{anotações-ouro}, ou \textit{conjunto-ouro} (\textit{gold set}): um conjunto de respostas do assistente cujos vereditos corretos foram atribuídos manualmente, montado separadamente para cada funcionalidade (contextualização, relembrar, definição e paráfrase). Nesse conjunto, a cobertura dos tipos de afirmação importa mais que o volume --- mais vale incluir um caso de cada situação difícil do que muitos exemplos fáceis. Por isso ele contém erros plantados deliberadamente: definições de sentido trocado, paráfrases que deslocam o significado, contextualizações sem suporte no texto e fatos de mundo não confirmáveis. Sem erros conhecidos seria impossível medir se o validador os detecta; plantá-los é o que torna a detecção mensurável.

Com esse conjunto, a calibração ocorre em três passos. Primeiro, rodamos o validador em cada item e marcamos cada veredito como certo ou errado, comparando-o ao rótulo humano; a confiança verbalizada do modelo é tomada como o \textit{score}. Segundo, antes de fixar qualquer limiar, verificamos se esse \textit{score} é sequer utilizável --- isto é, se ele de fato atribui confiança mais alta aos vereditos certos do que aos errados. Essa é a propriedade que a área sob a curva ROC (AUC) quantifica (\cite{Fawcett2006ROC}): um valor de $0{,}5$ indica que o \textit{score} de confiança não discrimina nada, e nenhum limiar seria útil; quanto mais próximo de $1$, melhor a confiança distingue os vereditos certos dos errados. Terceiro, confirmada a discriminação, lemos o limiar na curva de risco--cobertura: percorrendo os valores possíveis de $\tau$, observamos para cada um a fração de respostas em que o validador assume o veredito (cobertura) e a acurácia entre essas respostas, e escolhemos o $\tau$ mais permissivo que ainda respeita a acurácia desejada --- aquele que admite o máximo de respostas sem ultrapassar o nível de erro aceitável.

Duas ressalvas delimitam o alcance desse procedimento e valem ser comentadas aqui. A primeira é o tamanho: nosso \textit{conjunto-ouro} é pequeno e anotado por uma única pessoa, de modo que as métricas e os limiares obtidos são ilustrativos do método, não valores de produção --- um sistema real exigiria um conjunto bem maior, com vários anotadores. A segunda é o teto realista de desempenho: em afirmações genuinamente ambíguas, anotadores humanos discordam entre si, e não se espera que o validador supere essa concordância humano--humano. A medida dessa concordância é, portanto, o limite superior honesto contra o qual o desempenho do validador deve ser lido, e não a perfeição.


\subsubsection*{Discussão}
\label{subsubsec:discussao-validacao}

Vale recuperar o fio que percorre toda esta seção. O validador nasce de uma recusa em tratar ``correto'' como uma propriedade única: ao decompor cada resposta por tipo de afirmação e checar cada uma contra a fonte de evidência que lhe é própria, ele troca um juízo global e frágil por uma sequência de checagens menores, cada uma mais simples de fazer bem. Sobre essa base assenta o que consideramos o resultado mais importante do ponto de vista de IA confiável: a separação entre o veredito e a confiança que se pode depositar nele. É essa separação que dá sentido ao estado \textit{Com ressalvas} --- a admissão explícita de que uma resposta pode ser plausível e, ainda assim, não verificável --- e que distingue o nosso sistema de um assistente que apenas afirma, sem nunca sinalizar o quanto se deve acreditar no que afirma.

Igualmente central é o fato de que essa confiança não é assumida, mas \textit{conquistada} sobre dados: a confiança verbalizada do modelo, sabidamente superconfiante, só vira um sinal utilizável depois de passar pelo teste de discriminação e pela calibração do limiar no conjunto-ouro. Esse é o sentido prático de um nível de maturidade TRL 3 --- não a promessa de um validador infalível, mas um mecanismo cujo comportamento foi medido e caracterizado, com a honestidade de reportar onde ele acerta e de reconhecer, desde já, que o conjunto sobre o qual foi calibrado é pequeno demais para sustentar garantias de produção. Mais do que entregar respostas verificadas, o validador torna visível ao leitor o próprio ato de verificar --- e é nessa transparência, e não numa suposta certeza, que reside a sua contribuição para a confiabilidade do assistente.
